\documentclass{letter}
\usepackage[a4paper,
left=2cm,
right=2cm,
top=2.5cm,
bottom=2.5cm]{geometry}

\begin{document}



\large{\textbf{Nyilatkozat}}
\\

Ságodi Zoltán ,,Exploring Large Language Models’ Capabilities in Software Engineering'' című PhD disszertációjában a következő ered\-mé\-nyek\-ben Ságodi Zoltán hozzájárulása volt a meghatározó:

\begin{enumerate}
	\item Tézis:\\
		\textbf{Zoltán Ságodi}, István Siket, and Rudolf Ferenc
		\\ {Methodology for Code Synthesis Evaluation of LLMs Presented by a Case Study of ChatGPT and Copilot}.
		\\ \emph{IEEE Access}, VOL(12), 72303-72316, 2024.

		\begin{itemize}
			\item[I/1.] Kapcsolódó módszertanok összegyűjtése, és azok hi\-á\-nyos\-sá\-ga\-i\-nak meghatározása
		
			\item[I/2.] Összehasonlítási módszertan kidolgozása.
			
			\item[I/3.] Értékelési teszthalmaz kiválasztása.
			
			\item[I/4.] A nyelvi modellek kiértékelése a teszthalmaz segítségével.
			
			\item[I/5.] Szintetizált programok minőségi elemzése.
			
			\item[I/6.] Az emberi értékelési folyamat megtervezése.
			
			\item[I/7.] Az emberi értékelési keretrendszer megvalósítása.
		\end{itemize}
\vspace{1em}
	\item Tézis:\\
	\textbf{Zoltán Ságodi}, István Kolláth, Péter Heged\H{u}s, and Rudolf Ferenc\\
	{A Program Synthesis Dataset for LLM Temperature Analysis}.\\
	\emph{IEEE Access}, VOL(13), 184022-184029, 2025.
	
		\begin{itemize}
			\item[II/1.] Kutatás keretének megtervezése.
			
			\item[II/2.] Modellválasztó keretrendszer elkészítése.
			
			\item[II/3.] Modell kiértékelő keretrendszer elkészítése.
			
			\item[II/4.] Az adat kiértékelése a temperature hiperparaméterre fókuszálva.
		\end{itemize}
\vspace{1em}
	\item Tézis:\\
		\textbf{Zoltán Ságodi}, Gábor Antal, Bence Bogenfürst, Martin Isztin, Péter Heged\H{u}s, and Rudolf Ferenc\\
		{Reality Check: Assessing GPT-4 in Fixing Real-World Software Vulnerabilities}.\\
		In \emph{Proceedings of the 28th International Conference on Evaluation and Assessment in Software Engineering (EASE '24)}, Association for Computing Machinery, 252–261, 2024.
	
		\begin{itemize}
			\item[III/1.] Promptolási technikák gyűjtése.
			
			\item[III/2.] Promptok és modellek kiértékelésének tervezése.
			
			\item[III/3.] Részvétel a generált kódok kézi kiértékelésében.
			
			\item[III/4.] Részvétel a generált szöveges válaszok kézi kiértékelésében.
		\end{itemize}
	
\end{enumerate}
\newpage
Ezek az eredmények Ságodi Zoltán PhD disszertációján kívül más tudományos fokozat
megszerzésére nem használhatók fel.
\vspace{2em}

------------------------------------------------------------------------------------------------------------------\\
dátum, jelölt aláírása, témavezetők aláírása
\vspace{0.5em}

Az Informatika Doktori Iskola vezetője kijelenti, hogy jelen nyilatkozatot minden
társszerzőhöz eljuttatta, és azzal szemben egyetlen társszerző sem emelt kifogást.
\vspace{2em}

---------------------------------------------------------------\\
dátum, DI vezető aláírás

\end{document}